\chapter{فرمول‌ها و فرض‌هایی که پشت آن‌ها قرار دارند}
\label{app:information-formulas}

در متن فصل، چند رابطه‌ی ریاضی را برای سنجیدن اطلاعات، توصیف کانال و توضیح نمونه‌برداری به کار بردیم. این روابط در نگاه اول ممکن است شبیه فرمول‌هایی مستقل از یکدیگر به نظر برسند؛ اما هرکدام برای یک مدل مشخص ساخته شده‌اند. بنابراین پیش از آنکه از آن‌ها در مسئله‌ای تازه استفاده کنیم، باید بدانیم چه چیزی را اندازه می‌گیرند و چه فرض‌هایی پشت آن‌ها قرار دارد.

\section{تعداد حالت‌ها و مقدار اطلاعات}

اگر هر پیام از \(n\) جایگاه تشکیل شده باشد و در هر جایگاه بتوان یکی از \(s\) نماد را قرار داد، تعداد رشته‌های ممکن برابر است با:

\[
M=s^n.
\]

اگر همه‌ی این رشته‌ها به یک اندازه محتمل باشند، مقدار اطلاعات لازم برای تشخیص یکی از آن‌ها برابر خواهد بود با:

\[
I=\log_2 M=n\log_2 s.
\]

این رابطه به فرض برابری احتمال‌ها وابسته است. اگر بعضی پیام‌ها بسیار محتمل‌تر از بقیه باشند، شمارش ساده‌ی حالت‌ها تصویر کاملی از عدم‌قطعیت منبع به ما نمی‌دهد و باید احتمال هر پیام را نیز وارد محاسبه کنیم.

\section{خوداطلاعات و آنتروپی شانون}

برای یک نتیجه‌ی مشخص \(x_i\) با احتمال \(p_i\)، خوداطلاعات آن نتیجه به‌صورت زیر تعریف می‌شود:

\[
I(x_i)=-\log_2 p_i.
\]

این کمیت مقدار غافلگیریِ همان نتیجه را در مدل احتمالاتی ما توصیف می‌کند؛ نتیجه‌ای که احتمال کمتری دارد، خوداطلاعات بیشتری خواهد داشت. اگر نتیجه از پیش قطعی باشد، \(p_i=1\) و خوداطلاعات آن صفر است.

آنتروپی شانون میانگین خوداطلاعات نتیجه‌های ممکن است:

\[
H(X)=-\sum_i p_i\log_2 p_i.
\]

در این رابطه، \(X\) متغیر تصادفی، \(x_i\) یکی از نتیجه‌های ممکن و \(p_i\) احتمال همان نتیجه است. اگر لگاریتم بر پایه‌ی دو باشد، واحد آنتروپی «بیت» خواهد بود. آنتروپی به توزیع احتمال وابسته است، نه به معنای زبانی یا اهمیت انسانی پیام.

برای یک متغیر دودویی با احتمال \(p\) برای یکی از دو نتیجه، آنتروپی دودویی چنین نوشته می‌شود:

\[
H_2(p)=-p\log_2 p-(1-p)\log_2(1-p).
\]

در این فرمول، قرارداد \(0\log_2 0=0\) در نظر گرفته می‌شود؛ بنابراین رابطه در دو انتهای بازه نیز به شکل پیوسته تعریف‌پذیر است.

\section{ظرفیت کانال و تفاوت دو مدل}

برای یک کانال دودویی متقارنِ بی‌حافظه، اگر احتمال وارونه‌شدن هر نماد \(p\) باشد، ظرفیت بر حسب بیت در هر بار استفاده از کانال برابر است با:

\[
C=1-H_2(p).
\]

این رابطه را نباید بدون توجه به مدل کانال به هر مسیر ارتباطی تعمیم داد. فرض آن این است که کانال در هر استفاده یکی از دو نماد را دریافت می‌کند، خطا در هر دو جهت احتمال یکسان دارد و خطاهای استفاده‌های مختلف مستقل‌اند.

برای یک کانال پیوسته‌ی ایدئال با نویز سفید گاوسی، رابطه‌ی شانون ـ هارتلی به شکل زیر بیان می‌شود:

\[
C=B\log_2\left(1+\frac{S}{N}\right).
\]

در این رابطه \(B\) پهنای باند بر حسب هرتز و \(S/N\) نسبت توان سیگنال به توان نویز است. نسبت \(S/N\) باید به شکل خطی و نه بر حسب دسی‌بل در این فرمول قرار گیرد. خروجی \(C\) در این مدل، نرخ ظرفیت بر حسب بیت بر ثانیه است. این رابطه به مدل کانال و محدودیت توان وابسته است و فرمول عمومی همه‌ی کانال‌های واقعی نیست.

\section{پیچیدگی کولموگروف}

پیچیدگی کولموگروف یک رشته‌ی \(x\) نسبت به ماشین توصیف‌گر همگانی \(U\)، طول کوتاه‌ترین برنامه‌ای است که آن رشته را تولید می‌کند:

\[
K_U(x)=\min_{p:U(p)=x}|p|.
\]

در این رابطه \(p\) یک برنامه، \(U(p)\) خروجی اجرای آن برنامه و \(|p|\) طول توصیف برنامه بر حسب نمادهای زبان برنامه است. این کمیت با آنتروپی شانون یکی نیست: آنتروپی درباره‌ی توزیع احتمال یک منبع است، درحالی‌که پیچیدگی کولموگروف به کوتاه‌ترین توصیف یک رشته‌ی مشخص نگاه می‌کند. مقدار دقیق پیچیدگی به ماشین مرجع وابسته است، هرچند برای ماشین‌های همگانی مناسب اختلاف آن‌ها حداکثر تا یک مقدار ثابت محدود می‌شود.

\section{پاک‌کردن اطلاعات و اصل لاندائر}

برای پاک‌کردن منطقی یک بیتِ ناشناخته در دمای \(T\)، حد پایین انرژی گرماییِ دفع‌شده در محیط، در شرایط اصل لاندائر، برابر است با:

\[
Q_{\min}=k_{\mathrm B}T\ln 2.
\]

در این رابطه \(k_{\mathrm B}\) ثابت بولتزمن و \(T\) دمای مطلق بر حسب کلوین است. این رابطه درباره‌ی یک عملیات منطقیِ برگشت‌ناپذیر مانند پاک‌کردن اطلاعات است؛ نباید آن را به این ادعا تبدیل کرد که هر محاسبه یا هر جابه‌جایی اطلاعات الزاماً همین مقدار انرژی مصرف می‌کند. مقدار واقعی انرژی در یک دستگاه به طراحی، سرعت، دما و سازوکار فیزیکی آن وابسته است.

\section{نمونه‌برداری و شرط بازسازی}

اگر از یک سیگنال در هر ثانیه \(f_s\) نمونه بگیریم، فاصله‌ی زمانی میان نمونه‌های متوالی برابر است با:

\[
T_s=\frac{1}{f_s}.
\]

برای سیگنالی که به‌طور ایده‌آل به فرکانس‌های پایین‌تر از \(f_{\max}\) محدود شده باشد، شرط نمونه‌برداری نایکوئیست ـ شانون برای بازسازی ایده‌آل چنین بیان می‌شود:

\[
f_s>2f_{\max}.
\]

این شرط به محدودبودن باند سیگنال و فرض‌های بازسازی ایده‌آل وابسته است. در سامانه‌های واقعی، فیلترکردن پیش از نمونه‌برداری، محدودیت قطعات و حاشیه‌ی اطمینان نیز اهمیت دارند. صرفاً افزایش نرخ نمونه‌برداری، اگر سیگنال به‌درستی آماده نشده باشد، همه‌ی خطاها را از بین نمی‌برد.

\chapter{کمیت‌ها و واحدهای فیزیکی}
\label{app:physical-quantities}

در بحث الکتریسیته، نوشتن یک عدد به‌تنهایی کافی نیست. باید بدانیم آن عدد کمیتِ چه چیزی است و با چه واحدی بیان شده است. چند واحد زیر در فصل به کار رفتند یا برای خواندن رابطه‌های آن ضروری‌اند.

\begin{description}
  \item[\(\mathrm{bit}\)] واحد مقدار اطلاعات در صورت استفاده از لگاریتم بر پایه‌ی دو. بیت در اینجا یک واحد اطلاعاتی است، نه یک واحد پایه‌ی دستگاه بین‌المللی یکاها.
  \item[\(\mathrm{Hz}\)] هرتز؛ واحد بسامد و برابر با یک چرخه در ثانیه، یعنی \(\mathrm{s}^{-1}\). نرخ نمونه‌برداری نیز با نمونه بر ثانیه بیان می‌شود.
  \item[\(\mathrm{V}\)] ولت؛ واحد اختلاف پتانسیل الکتریکی. یک ولت برابر یک ژول بر کولن است:
  \[
  1\,\mathrm{V}=1\,\frac{\mathrm{J}}{\mathrm{C}}.
  \]
  \item[\(\mathrm{C}\)] کولن؛ واحد بار الکتریکی. در فرمول‌های اطلاعاتی، \(C\) بزرگ همچنین برای ظرفیت کانال به کار می‌رود؛ زمینه‌ی رابطه مشخص می‌کند کدام معنا منظور است.
  \item[\(\mathrm{J}\)] ژول؛ واحد انرژی و کار.
  \item[\(\mathrm{K}\)] کلوین؛ واحد دمای ترمودینامیکی.
  \item[\(\mathrm{s}\)] ثانیه؛ واحد زمان.
\end{description}

\section{پیشوندهای ده‌دهی پرکاربرد}

پیشوندها اندازه‌ی یکای اصلی را در توان‌های ده تغییر می‌دهند. برای نمونه:

\[
\begin{array}{c|c|c}
\text{پیشوند} & \text{نماد} & \text{ضریب}\\
\hline
\text{کیلو} & \mathrm{k} & 10^3\\
\text{مگا} & \mathrm{M} & 10^6\\
\text{گیگا} & \mathrm{G} & 10^9\\
\text{میلی} & \mathrm{m} & 10^{-3}\\
\text{میکرو} & \mathrm{\mu} & 10^{-6}\\
\text{نانو} & \mathrm{n} & 10^{-9}
\end{array}
\]

برای مثال، \(48\,\mathrm{kHz}\) یعنی چهل‌وهشت هزار هرتز. حروف بزرگ و کوچک در نماد پیشوندها اهمیت دارند؛ \(\mathrm{M}\) و \(\mathrm{m}\) یک معنا ندارند.

\section{اختلاف پتانسیل و انرژی}

اگر بار \(q\) میان دو نقطه جابه‌جا شود و تغییر انرژی پتانسیل الکتریکی آن \(\Delta U\) باشد، اختلاف پتانسیل میان آن دو نقطه چنین تعریف می‌شود:

\[
\Delta V=\frac{\Delta U}{q},
\qquad
\Delta U=q\Delta V.
\]

بنابراین ولتاژ همان انرژی نیست؛ انرژی پتانسیل الکتریکی به مقدار بار نیز وابسته است. واحدها با یکدیگر سازگارند: ژول تقسیم بر کولن، ولت به دست می‌دهد.
